\section*{第 \ref{ch:b1:plant-water-minerals} 章——植物水分与矿质营养}

\begin{solution}{exo:b1:plant-water-minerals:1}
\omterm{def:b1:plant-water-minerals:pathways}{质外体}：由连续的\omterm{def:b1:eukaryotic-cell:plantcell}{细胞壁}与空隙组成，物质
无需穿过膜即可通过。\omterm{def:b1:plant-water-minerals:pathways}{共质体}：各\omterm{def:b1:cell-unit-of-life:cell}{细胞}相连的\omterm{def:b1:cell-unit-of-life:cell}{细胞质}，
经一层膜进入，并经\omterm{def:b1:eukaryotic-cell:plantcell}{胞间连丝}延续。\omterm{def:b1:plant-water-minerals:pathways}{凯氏带}
阻断\omterm{prop:b1:flowering-plant-organization:stemroot}{内皮层}处的\omterm{def:b1:plant-water-minerals:pathways}{质外体}途径，迫使
一切物质穿过膜，并防止溶液从
\omterm{def:b1:flowering-plant-organization:tissues}{木质部}回漏。
\end{solution}

\begin{solution}{exo:b1:plant-water-minerals:2}
N（\omterm{def:b1:proteins:peptide}{蛋白质}、核酸）、P（\omterm{def:b1:nucleic-acids:nucleotide}{核苷酸}、膜）、K（渗透剂、
\omterm{def:b1:flowering-plant-organization:tissues}{气孔}）、Ca（\omterm{def:b1:eukaryotic-cell:plantcell}{细胞壁}、信号传导）、Mg（\omterm{def:b1:photosynthesis:pigments}{叶绿素}）、
S（\omterm{def:b1:proteins:aminoacid}{氨基酸}：半胱氨酸、甲硫氨酸）。
\end{solution}

\begin{solution}{exo:b1:plant-water-minerals:3}
土壤 \qty{-0.2}{}、根\omterm{def:b1:flowering-plant-organization:tissues}{木质部} \qty{-0.6}{}、叶\omterm{def:b1:cell-unit-of-life:cell}{细胞} \qty{-1.3}{MPa}：
水分自土壤进入根，再自根\omterm{def:b1:flowering-plant-organization:tissues}{木质部}上升至叶——总是
朝着更负的数值方向流动。
\end{solution}

\begin{solution}{exo:b1:plant-water-minerals:4}
氮在\omterm{def:b1:flowering-plant-organization:tissues}{韧皮部}中可移动：植物将其从老叶中调出，
供给幼叶，故老叶先变黄。铁在体内不能
再利用：新长出的幼叶因缺铁而
变黄。
\end{solution}

\begin{solution}{exo:b1:plant-water-minerals:5}
\omterm{def:b1:cell-unit-of-life:cell}{细胞} $\Psi = -0.9 + 0.5 = \qty{-0.4}{MPa}$：土壤（$-0.2$）$>$ \omterm{def:b1:cell-unit-of-life:cell}{细胞}
（$-0.4$）$>$ \omterm{def:b1:flowering-plant-organization:tissues}{木质部}（$-0.6$），因此水分由土壤 $\to$ \omterm{def:b1:cell-unit-of-life:cell}{细胞} $\to$
\omterm{def:b1:flowering-plant-organization:tissues}{木质部}。当\omterm{def:b1:cell-unit-of-life:cell}{细胞}达到 \qty{-0.2}{MPa} 时吸水停止，即
$\Psi_p = -0.2 + 0.9 = \qty{0.7}{MPa}$。
\end{solution}

\begin{solution}{exo:b1:plant-water-minerals:6}
$\Psi = 137\ln(\text{RH})$：\qty{-14}{MPa}，\qty{-95}{MPa}，
\qty{-315}{MPa}。即使潮湿的空气，按此意义也比
萎蔫点时的土壤干十倍：叶片总是向空气失水，
唯一的问题是失水的速度。
\end{solution}

\begin{solution}{exo:b1:plant-water-minerals:7}
$E_K = 25.7\ln(0.1/100) = \qty{-178}{mV}$：\omterm{def:b1:cell-unit-of-life:cell}{细胞}在 \qty{-150}{mV}
时的电位比所需值略低（负得不够），因此千倍的积累
还需要略多于通道的作用——但比值为 100 时即处于
平衡：通道几乎完成了全部工作。硝酸根：$E = -25.7\ln(0.5/5)
= \qty{+59}{mV}$；\omterm{def:b1:cell-unit-of-life:cell}{细胞}位于 \qty{-150}{}，与之相差 \qty{209}{mV}，且
方向相反：阴离子必须借助质子同向转运被泵入。
\end{solution}

\begin{solution}{exo:b1:plant-water-minerals:8}
夜间无蒸腾，向\omterm{def:b1:flowering-plant-organization:tissues}{木质部}装载离子会吸引
水分进入，而约零点几兆帕的\omterm{prop:b1:plant-water-minerals:rootpressure}{根压}将汁液
最多推升几米——足以使草叶尖端溢出液滴，
却不能到达树顶。\omterm{def:b1:flowering-plant-organization:tissues}{气孔}张开时，蒸腾使\omterm{def:b1:flowering-plant-organization:tissues}{木质部}
承受张力，
\omterm{prop:b1:plant-water-minerals:rootpressure}{根压}随之消失。
\end{solution}

\begin{solution}{exo:b1:plant-water-minerals:9}
缺钼时\omterm{prop:b1:plant-water-minerals:nitrate}{硝酸还原酶}不能工作：以硝酸盐为氮源时
植物无法合成\omterm{def:b1:proteins:aminoacid}{氨基酸}，表现出缺氮，尽管硝酸盐充足。
以铵为氮源时植物生长正常，因为铵
无需还原即可进入\omterm{def:b1:proteins:aminoacid}{氨基酸}。这一差异将该元素
定位于一种酶。
\end{solution}

\begin{solution}{exo:b1:plant-water-minerals:10}
\qty{30}{mg} 氮 $= \qty{2.14}{mmol}$；$\times 8 = \qty{17}{mmol}$ 的
电子，\qty{8.6}{mmol} 的 \omterm{def:b1:biosyntheses-integration:ppp}{NADPH} 当量。叶片：$0.05\times
4\times 10^{-6}\times 43\,200 = \qty{8.6}{mmol}$ 的 $\mathrm{CO_2}$，
\qty{35}{mmol} 电子：硝酸盐还原所用的电子
数约为碳固定的一半——占比不小，
这正是植物在光下进行硝酸盐还原的原因，
也是速生叶片大量消耗还原剂的缘故。
\end{solution}

\begin{solution}{exo:b1:plant-water-minerals:11}
\qty{2.14}{mmol} 氮 $= \qty{1.07}{mmol}$ 的 $\mathrm{N_2}$：需
\omterm{def:b1:water-small-molecules:atp}{ATP} $16\times 1.07 = \qty{17}{mmol}$，即 \qty{0.57}{mmol} 葡萄糖，
\qty{0.10}{g}；另每个 $\mathrm{N_2}$ 还需 8 个电子（相当于
\qty{0.05}{g} 葡萄糖）。按 \qty{6}{g/g} 的规则，每天
约 \qty{0.18}{g}：\omterm{def:b1:plant-water-minerals:nodules}{固氮}酶本身占大部分，其余用于构建和
维持\omterm{def:b1:plant-water-minerals:nodules}{根瘤}、其中的豆血红蛋白，以及细菌自身
为维持低氧而进行的呼吸。
\end{solution}

\begin{solution}{exo:b1:plant-water-minerals:12}
根自身吸收水分，以及土壤水分带到\omterm{def:b1:flowering-plant-organization:organs}{根毛}处的可移动
离子（硝酸根、钾），仅在\omterm{prop:b1:plant-water-minerals:uptake}{质子泵}上消耗
\omterm{def:b1:water-small-molecules:atp}{ATP}。对于不会自行到达根部的磷酸盐，根以糖
支付真菌报酬，借其扩展吸收范围；对于缺乏硝酸盐土壤中的氮，
则以糖向细菌支付报酬，令其固定氮。费用占
光合产物的十分之一或更多，而植物在不再需要
服务时即停止支付：磷丰富的土壤中，\omterm{def:b1:plant-water-minerals:mycorrhiza}{菌根}减少；硝酸盐
丰富的土壤中，豆科植物\omterm{def:b1:plant-water-minerals:nodules}{根瘤}减少——承包者
只在比自行完成工作更便宜时才被保留。
\end{solution}

\begin{solution}{pb:b1:plant-water-minerals:1}
\textbf{1.} $-0.9 + 0.5 = \qty{-0.4}{MPa}$。
\textbf{2.} $-0.2 > -0.4 > -0.6 > -1.3$：逐级降低，水分流动：
土壤 $\to$ 皮层 $\to$ \omterm{def:b1:flowering-plant-organization:tissues}{木质部} $\to$ 叶。
\textbf{3.} $137\ln 0.5 = \qty{-95}{MPa}$。
\textbf{4.} 土壤到叶 \qty{1.1}{MPa}；叶到空气 \qty{94}{MPa}：
总量的 \qty{99}{\%} 来自最后一步。
\textbf{5.} 土壤的 $-0.9$ 低于皮层\omterm{def:b1:cell-unit-of-life:cell}{细胞}的 $-0.4$：水分
离开\omterm{def:b1:cell-unit-of-life:cell}{细胞}，直到\omterm{def:b1:cell-unit-of-life:cell}{细胞}降至 $-0.9$，即其\omterm{def:b1:eukaryotic-cell:plantcell}{膨压}降为
零；吸水停止，直到\omterm{def:b1:cell-unit-of-life:cell}{细胞}通过
积累溶质降低其 $\Psi_s$。
\textbf{6.} 此时土壤低于叶片的 $-1.3$：水分流出
植物体；叶片失去\omterm{def:b1:eukaryotic-cell:plantcell}{膨压}而萎蔫。浇水使
土壤的 $\Psi$ 高于根的，水流恢复；叶片的
\omterm{def:b1:eukaryotic-cell:plantcell}{细胞壁}完好，可以重新充盈。
\textbf{7.} $2\times 10^{-7}\times 0.2\times(-0.2 - (-0.6)) =
\qty{1.6e-8}{m^3/s}$。
\textbf{8.} $1.6\times 10^{-8}\times 86\,400 = \qty{1.4e-3}{m^3} =
\qty{1.4}{L}$ 每天。
\textbf{9.} \qty{1400}{g}，分摊于 \qty{0.5}{m^2} 面积、
\qty{12}{h} 日照（大部分时间）内：约每平方米每小时 \qty{230}{g}。
\textbf{10.} $230/18 = \qty{12.8}{mol}$ 每平方米每小时，\qty{3.5}{mmol} 每平方米每秒——即每固定 \qty{4}{\micro mol} 的 $\mathrm{CO_2}$，约散失一千个水分子。
\textbf{11.} 差值 $-0.2 - (-0.4) = 0.2$：通量减半，
每天 \qty{0.7}{L}。\omterm{def:b1:flowering-plant-organization:tissues}{气孔}关闭减少了失水，张力
得以缓解，整个\omterm{def:b1:membranes-transport:osmosis}{水势}阶梯也已抬高：植物在空气最干燥时
以牺牲碳增益换取水分。
\textbf{12.} 水分沿自身的\omterm{def:b1:membranes-transport:osmosis}{水势}梯度下行，该梯度由
叶片蒸腾所造成，因此根的吸水是被动的；硝酸根必须
逆浓度梯度和电位梯度被浓缩，
这只有与\omterm{prop:b1:plant-water-minerals:uptake}{质子泵}偶联才能支付。
\textbf{13.} $\qty{1.4}{L}\times\qty{2}{mmol/L} = \qty{2.8}{mmol}$，
\qty{39}{mg} 氮。
\textbf{14.} 是的，尚余约三分之一；多余部分贮存于
\omterm{def:b1:eukaryotic-cell:plantcell}{液泡}中（以硝酸根形式）或在\omterm{prop:b1:flowering-plant-organization:stemroot}{内皮层}被排除。
\textbf{15.} $\Delta G = RT\ln(5/2) + zFV = 2.48\times 0.92 + (-1)(96.5)(-0.15)
= 2.3 + 14.5 = \qty{16.8}{kJ/mol}$。
\textbf{16.} 每个质子：$RT\ln(10^{-7.5}/10^{-5.5}) + F(-0.15) =
-11.4 - 14.5 = \qty{-25.9}{kJ}$；两个质子：\qty{-52}{kJ}，为
所需的三倍。
\textbf{17.} $\qty{2.14}{mmol}\times 8 = \qty{17}{mmol}$ 电子。
\textbf{18.} $0.5\times 4\times 10^{-6}\times 43\,200 = \qty{86}{mmol}$ 的
$\mathrm{CO_2}$，\qty{346}{mmol} 电子：\qty{5}{\%} 用于
硝酸盐还原。
\textbf{19.} 这样可节省每个氮的 8 个电子；但
每吸收一个阳离子便释放一个质子，使根周围土壤酸化，
\omterm{def:b1:cell-unit-of-life:cell}{细胞}必须向外输出酸；且铵若积累便有毒，
因此须在根中立即同化。
\textbf{20.} \qty{1.07}{mmol} 的 $\mathrm{N_2}$ $\times 16 =
\qty{17}{mmol}$ 的 \omterm{def:b1:water-small-molecules:atp}{ATP}。
\textbf{21.} $6\times 0.03 = \qty{0.18}{g}$ 的糖，占
\qty{3}{g}。若硝酸盐浓度为 \qty{2}{mmol/L}，且计入其还原
成本（占电子的 \qty{5}{\%}），\omterm{def:b1:plant-water-minerals:nodules}{固氮}是一笔
不划算的交易，豆科植物形成的\omterm{def:b1:plant-water-minerals:nodules}{根瘤}很少；若无硝酸盐，
\omterm{def:b1:plant-water-minerals:nodules}{固氮}便是唯一途径，且 \qty{6}{\%} 的产出
相对整株植物的\omterm{def:b1:proteins:peptide}{蛋白质}而言代价低廉。
\textbf{22.} $1.4\times 2\times 10^{-6} = \qty{2.8}{\micro mol}$，
\qty{0.09}{mg} 磷——仅为需求量的四十分之一；植物
须依靠扩散至\omterm{def:b1:flowering-plant-organization:organs}{根毛}，而\omterm{def:b1:flowering-plant-organization:organs}{根毛}很快耗尽其周围的
磷，更主要的是依赖其\omterm{def:b1:plant-water-minerals:mycorrhiza}{菌根}真菌。
\textbf{23.} \omterm{def:b1:plant-water-minerals:pathways}{凯氏带}阻断\omterm{def:b1:plant-water-minerals:pathways}{质外体}，因此硝酸根只能
经\omterm{prop:b1:flowering-plant-organization:stemroot}{内皮层}的膜进入\omterm{def:b1:flowering-plant-organization:tissues}{木质部}，这些膜将其向内泵送；
同时，\omterm{def:b1:plant-water-minerals:pathways}{凯氏带}防止浓缩的\omterm{def:b1:flowering-plant-organization:tissues}{木质部}汁液沿\omterm{def:b1:eukaryotic-cell:plantcell}{细胞壁}
向外渗漏，从而维持该梯度。
\textbf{24.} 其\omterm{def:b1:flowering-plant-organization:tissues}{木质部}汁液将类似土壤溶液——稀薄、
未经筛选，含有钠及其他不需要的离子——且无法
保持浓缩；在盐渍土中，钠将自由进入叶片，
使植物中毒，而正常植物在
\omterm{prop:b1:flowering-plant-organization:stemroot}{内皮层}处将其排除。
\textbf{25.} 每天约 \qty{1.4}{L} 水分通过 \qty{0.2}{m^2}
根系，携带 \qty{39}{mg} 氮，而需求量为
\qty{30}{}，其还原约占叶片
光合电子的 \qty{5}{\%}。
\end{solution}
